\section*{S2 Proofs}
\label{supp:proofs}

This section proves the results of Section~3 of the paper: Proposition~\ref{prop:clean} and the guarantee for a vote within one stage (S2.0, S2.1), Proposition~\ref{thm:step} and its general version (S2.2), and Proposition~\ref{prop:recover-value} together with the result on the local peak test (S2.3). Each result is for a fixed input or fixed parameters. None covers the passage from one stage of ESF to the next, the weighted vote, or the search for the candidate line of the recovery step; the remarks that follow each proof say what it covers and what it does not.

\subsection*{S2.0 Proposition~\ref{prop:clean}: precise statement and proof}

Let $\mathcal F_{\ell-1}$ be the flagged set after stage $\ell-1$, with $\mathcal F_0$ the
covariate screen, $E_\ell=\{1,\dots,n\}\setminus\mathcal F_{\ell-1}$ the set from which stage
$\ell$ draws, $N_\ell=|E_\ell|$, $\varepsilon_\ell=|E_\ell\cap O|/N_\ell$, and $\mathcal H_{\ell-1}$ the data
together with everything drawn before stage $\ell$. A replicate of stage $\ell$ is obtained
by drawing uniform subsamples of size $m$ from $E_\ell$ until an admissible one is found; let
$q^{c}_\ell$ be the fraction of admissible subsamples among the clean subsamples of $E_\ell$,
and $q_\ell>0$ the fraction among all subsamples of $E_\ell$.

\emph{Proposition~\ref{prop:clean} (precise statement).}
Conditionally on $\mathcal H_{\ell-1}$, the subsamples of stage $\ell$ are independent, and each
is uniform on the admissible subsamples of size $m$ of $E_\ell$, however many units two of them
share. Let $p_\ell=\binom{N_\ell-|E_\ell\cap O|}{m}\big/\binom{N_\ell}{m}$ be the probability that a
uniform subsample of $E_\ell$ is clean, and $a_\ell=q^{c}_\ell/q_\ell$. Conditionally on
$\mathcal H_{\ell-1}$, each replicate of stage $\ell$ is clean with probability
$\pi_\ell=p_\ell a_\ell$, and the stage contains a clean replicate with probability
$1-(1-p_\ell a_\ell)^{B/L}$. Moreover $a_\ell\ge q^{c}_\ell$ and, when
$|E_\ell\cap O|\le N_\ell-m$,
$\{1-\varepsilon_\ell N_\ell/(N_\ell-m+1)\}^{m}\le p_\ell\le(1-\varepsilon_\ell)^{m}$.

\emph{The first statement.} Fix a replicate of stage $\ell$. Conditionally on
$\mathcal H_{\ell-1}$ its proposals are independent uniform subsamples of size $m$ of $E_\ell$,
because $E_\ell$ is determined by $\mathcal H_{\ell-1}$ and the draws use fresh randomness. (The
implementation draws $m$ units without replacement with equal weights on $E_\ell$ and zero
weight elsewhere, which gives a uniform subsample of $E_\ell$.) Whether a proposal is admissible
is determined by the proposal and the data. The first admissible proposal in a sequence of
independent proposals has the law of one proposal conditioned on admissibility, which is the
uniform law on the admissible subsamples; the sequence stops almost surely because $q_\ell>0$.
Different replicates use disjoint randomness, so they are conditionally independent. \qed

\emph{The rest.} By the first statement, a replicate of stage $\ell$ is clean with
conditional probability equal to the number of admissible clean subsamples divided by the
number of admissible subsamples. Dividing both by $\binom{N_\ell}{m}$ gives
$\pi_\ell=p_\ell q^{c}_\ell/q_\ell=p_\ell a_\ell$. The replicates of the stage are
conditionally independent, which gives the probability that at least one is clean. Next,
$a_\ell\ge q^{c}_\ell$ because $q_\ell\le1$. Finally
$p_\ell=\prod_{j=0}^{m-1}(N_\ell-|E_\ell\cap O|-j)/(N_\ell-j)$; each factor equals
$1-|E_\ell\cap O|/(N_\ell-j)$, which lies between $1-|E_\ell\cap O|/(N_\ell-m+1)$ and
$1-|E_\ell\cap O|/N_\ell=1-\varepsilon_\ell$. \qed

\emph{Remarks.} The proposition is exact and elementary. What matters for stage $\ell$ is the contamination
among the units still eligible: the flagged set lowers $\varepsilon_\ell$ when it contains
outliers, and the probability of a clean replicate rises geometrically in $m$ as it does so;
the proposition does not say that earlier stages lower $\varepsilon_\ell$. It does say where
ESF breaks down: the first stage draws from the whole sample, less the covariate screen, and
with $m=8$ its probability of containing a clean replicate is about $0.84$ at
$\varepsilon=0.20$, $0.45$ at $0.30$ and $0.28$ at $0.35$, and with $m=12$ about $0.51$,
$0.13$ and $0.06$ (taking $a_\ell=1$ and $p_\ell\approx(1-\varepsilon)^{m}$); if it finds no
clean replicate, its flags need not remove the outliers, and the breakdown observed from
$25\%$ to $30\%$ of outliers in some designs is of this order.

\subsection*{S2.1 A vote within one stage}

\emph{A vote within one stage.} The same conditional independence gives a guarantee for a vote within one stage. Suppose a
clean replicate gives unit $i$ its correct label with probability at least $\bar q$, where the
label refers to the least-squares partition of the clean eligible units. If
$\pi_\ell\bar q>\tfrac12$, a plurality vote of the replicates of stage $\ell$ labels every such
unit correctly with conditional probability at least $1-n\exp\{-2(B/L)(\pi_\ell\bar q-\tfrac12)^{2}\}$,
whatever the contaminated replicates do (proved below). Since
$\pi_\ell\approx(1-\varepsilon_\ell)^{m}$, this requires
\begin{equation}
m\;<\;\frac{\log(2\bar q)}{-\log(1-\varepsilon_\ell)}\;\approx\;\frac{\log(2\bar q)}{\varepsilon_\ell},
\label{eq:window}
\end{equation}
while admissibility and the accuracy of a clean replicate require $m$ of order at least
$(d+1)/\pi_{\min}$. The subsample size must lie between the two bounds. As the contamination among the eligible
units grows, the range of admissible $m$ shrinks; each stage of flagging lowers that
contamination and so widens the range again. The window can be empty: with $\bar q$ close to one, \eqref{eq:window} requires
$\varepsilon_\ell<1-2^{-1/m}$, which is $0.083$ for $m=8$ and $0.056$ for $m=12$, so at the
contamination levels of our experiments the guarantee applies only at stages where flagging has
already removed most of the outliers. This guarantee covers a single stage, not the weighted vote across stages that
Algorithm~\ref{alg:adaptive} uses; it motivates the design but says nothing about the final
output.

\emph{The vote within one stage.} Let $t_i$ be the label of unit $i$ in the least-squares
partition of the clean eligible units, labelled by the same pilot as the replicates, and
suppose a clean replicate gives unit $i$ the label $t_i$ with conditional probability at
least $\bar q$. Then a replicate of stage $\ell$ gives unit $i$ the label $t_i$ with
conditional probability $Q_i\ge\pi_\ell\bar q$, whatever the contaminated replicates do, since
the event that the replicate is clean and labels $i$ correctly has probability $\pi_\ell$ times
the conditional probability of a correct label given cleanness. By the first statement of Proposition~\ref{prop:clean}, the number of
the $B/L$ replicates that give unit $i$ the label $t_i$ is binomial with success probability
$Q_i$. If $Q_i\ge\tfrac12+\gamma$ with $\gamma>0$, the plurality differs from $t_i$ only if
at most half of the replicates vote for $t_i$, which by Hoeffding's inequality has
probability at most $\exp\{-2(B/L)\gamma^{2}\}$. A union bound over at most $n$ units, with
$\gamma=\pi_\ell\bar q-\tfrac12$, gives the statement above. The window
\eqref{eq:window} is $(1-\varepsilon_\ell)^{m}\bar q>\tfrac12$ solved for $m$, with $a_\ell=1$. \qed

\subsection*{S2.2 The final reweighting step}
\label{app:reweight}

Proposition~\ref{thm:step} of Section~\ref{sec:reweight} collects parts of
Theorems~\ref{thm:stepfixed} and~\ref{thm:steptrue} below.

The last two lines of Algorithm~\ref{alg:adaptive} apply a reweighting step to the
parameter produced by the concentration steps. We study the step as a map on its own,
applied to a fixed parameter $\theta=(\theta_1,\dots,\theta_K)\in\R^{Kd}$, and let the
sample size grow. The step, as implemented, is the following. For unit $i$ let
$z_i(\theta)=\argmin_k|y_i-x_i^{\top}\theta_k|$ be its nearest surface, the smallest index in
case of a tie, and $r_i(\theta)=\min_k|y_i-x_i^{\top}\theta_k|$ its absolute residual from
that surface. Let $n_k=\#\{i:z_i(\theta)=k\}$. The scale of group $k$ is
$\hat s_k=1.4826\,\hat m_k$, where $\hat m_k$ is a sample median of
$\{r_i(\theta):z_i(\theta)=k\}$, by which we mean any value between the two middle order
statistics of that set; the implementation takes the square root of the median of the
squared residuals, which lies between them. When $n_k<10$ the median over all $n$ residuals
is used instead. The flagged set and the flagged fraction are
\[
\mathcal F_n(\theta)=\{i:\ r_i(\theta)>c\,\hat s_{z_i(\theta)}\},\qquad
\hat\alpha_n(\theta)=|\mathcal F_n(\theta)|/n .
\]
The refitted parameter $\hat\theta_k$ is the least-squares fit of $y$ on $x$ over
$\{i:z_i(\theta)=k,\ i\notin\mathcal F_n(\theta)\}$ when this set has at least $d$ units, and
$\hat\theta_k=\theta_k$ otherwise.

The population counterparts are defined for a probability measure $P$ on
$\mathcal U=\R^{d}\times\R$ and a unit $u=(x,y)$:
\begin{align*}
h_\theta(u)&=\argmin_k|y-x^{\top}\theta_k|,\qquad
r_\theta(u)=\min_k|y-x^{\top}\theta_k|,\qquad
A_k(\theta)=\{u:h_\theta(u)=k\},\\
m_k(\theta)&=\inf\{t\ge0:\ P(r_\theta\le t\mid A_k(\theta))\ge\tfrac12\},\qquad
s_k(\theta)=1.4826\,m_k(\theta),\\
W_k(\theta)&=A_k(\theta)\cap\{r_\theta\le c\,s_k(\theta)\},\qquad
\alpha(\theta)=P\bigl(r_\theta>c\,s_{h_\theta}(\theta)\bigr)=1-\textstyle\sum_kP(W_k(\theta)),\\
M_k(\theta)&=\E\bigl[xx^{\top}\one_{W_k(\theta)}\bigr],\qquad
v_k(\theta)=\E\bigl[xy\,\one_{W_k(\theta)}\bigr],\qquad
\theta^{\infty}_k(\theta)=M_k(\theta)^{-1}v_k(\theta).
\end{align*}
Here $m_k(\theta)$ is the median of the residual $r_\theta$ among the units assigned to
surface $k$, $s_k(\theta)$ the corresponding scale, $W_k(\theta)$ the set of units assigned to
$k$ and not flagged, $\alpha(\theta)$ the population flagged fraction, and
$\theta^{\infty}(\theta)$ the population refit. The dependence on $\theta$ is dropped when
$\theta$ is fixed.

\begin{assumption}[Clean law]
\label{ass:clean}
$P=(1-\varepsilon)P_0+\varepsilon C$ with $\varepsilon\in[0,1)$. Under $P_0$ a unit is
generated by a label $z\in\{1,\dots,K\}$ with $\Pr(z=k)=\pi_k>0$, a covariate $x$ whose
conditional law given $z=k$ has $\E_0[\|x\|^{2}]<\infty$, and an error $e\sim\mathcal N(0,1)$
independent of $(x,z)$, through $y=x^{\top}\theta^{0}_z+\sigma_ze$ with $\sigma_k>0$. The
contaminating law $C$ is arbitrary.
\end{assumption}

\begin{assumption}[Contamination beyond the cut-off]
\label{ass:far}
At the given $\theta$, there is $c'>c$ such that
$C\bigl(r_\theta\ge c'\,s_{h_\theta}(\theta)\bigr)=1$: every contaminated unit lies beyond
$c'$ scales from every surface of $\theta$, the scale being that of the surface nearest to it.
\end{assumption}

The scale $s_k(\theta)$ in Assumption~\ref{ass:far} is computed under $P$, so the
assumption is a condition on $P$; it is the form of ``beyond the cut-off'' that refers to the
scale the step actually uses. In the designs of the study the gross outliers are shifted by $15$
to $30$ noise scales from their line (moderate outliers by $4$ to $8$), and uniform noise is
redrawn until it lies more than three noise scales from both true lines, so at $\theta^{0}$ the
assumption holds for all or nearly all contaminated units whenever $s_k(\theta^{0})$ is close to
$\sigma_k$; a shifted unit can still fall near another group's line, and the assumption then
fails for that unit.

\begin{assumption}[Regularity at $\theta$]
\label{ass:reg}
For every $k$: $P(A_k(\theta))>0$; the median of $r_\theta$ on $A_k(\theta)$ is positive and
unique, that is, $m_k(\theta)>0$, $P(r_\theta\le t\mid A_k)<\tfrac12$ for $t<m_k$ and
$>\tfrac12$ for $t>m_k$; and $M_k(\theta)$ is nonsingular.
\end{assumption}

\begin{lemma}[Plug-in cut-off]
\label{lem:plugin}
Let $u_1,u_2,\dots$ be independent with law $P$, let $k$ be fixed, let $\hat s_k$ be random
with $\hat s_k\to s_k$ almost surely, and let $g:\mathcal U\to\R^{p}$ be measurable with
$P(A_k,\,r_\theta=cs_k)=0$ and $\E[\|g\|\one\{A_k,\,r_\theta\le cs_k+\delta_0\}]<\infty$ for
some $\delta_0>0$. Then, almost surely,
\[
\frac1n\sum_{i=1}^{n}g(u_i)\one\{z_i(\theta)=k,\ r_i(\theta)\le c\hat s_k\}
\;\to\;\E\bigl[g\,\one_{W_k}\bigr].
\]
\end{lemma}

\begin{proof}
Fix $\delta\in(0,\delta_0)$. Almost surely $|c\hat s_k-cs_k|\le\delta$ for all large $n$, and
then the average on the left differs from
$n^{-1}\sum_ig(u_i)\one\{z_i=k,\,r_i\le cs_k\}$ by at most
$n^{-1}\sum_i\|g(u_i)\|\one\{z_i=k,\,cs_k-\delta<r_i\le cs_k+\delta\}$ in norm. By the strong
law of large numbers the second average converges to $\E[g\one_{W_k}]$ and the bound to
$\E[\|g\|\one\{A_k,\,cs_k-\delta<r_\theta\le cs_k+\delta\}]$, both integrable by hypothesis.
Letting $\delta\downarrow0$ along a sequence, the last expectation tends to
$\E[\|g\|\one\{A_k,\,r_\theta=cs_k\}]=0$ by dominated convergence.
\end{proof}

\begin{theorem}[One reweighting step at a fixed input]
\label{thm:stepfixed}
Let $u_1,\dots,u_n$ be independent with law $P$, fix $\theta$, and let
Assumptions~\ref{ass:clean}--\ref{ass:reg} hold at $\theta$. Then, almost surely as
$n\to\infty$:
\begin{enumerate}[label=(\roman*),leftmargin=2.2em,itemsep=2pt]
\item $\hat s_k\to s_k(\theta)$ for every $k$;
\item $\hat\alpha_n(\theta)\to\alpha(\theta)=\varepsilon+(1-\varepsilon)\kappa(\theta)$,
where $\kappa(\theta)=P_0\bigl(r_\theta>c\,s_{h_\theta}(\theta)\bigr)$ is the flagged
fraction among clean units;
\item $\hat\theta\to\theta^{\infty}(\theta)$, and
$\theta^{\infty}_k(\theta)=\E_0[xx^{\top}\one_{W_k}]^{-1}\E_0[xy\,\one_{W_k}]$: the limit of
the refit is an average over the clean law alone, and the contamination enters it only through
the scales $s_k(\theta)$ that define $W_k$.
\end{enumerate}
\end{theorem}

\begin{proof}
\emph{Step 1: no mass at the cut-off.} Under $P_0$, $y$ has a density conditionally on
$(x,z)$, so $P_0(|y-x^{\top}\theta_j|=t)=0$ for every $j$ and $t$, and
$P_0(r_\theta=t)\le\sum_jP_0(|y-x^{\top}\theta_j|=t)=0$. Under $C$,
Assumption~\ref{ass:far} gives $C(A_k,\,r_\theta<c's_k)=0$. Hence
$P(A_k,\,r_\theta=cs_k)=0$, and for $\delta_0=\tfrac12(c'-c)\min_ks_k$, which is positive by
Assumption~\ref{ass:reg}, and every $g\ge0$,
\begin{equation}
\E\bigl[g\,\one\{A_k,\,r_\theta\le cs_k+\delta_0\}\bigr]
=(1-\varepsilon)\E_0\bigl[g\,\one\{A_k,\,r_\theta\le cs_k+\delta_0\}\bigr].
\label{eq:adp-cleanonly}
\end{equation}
In particular $P(W_k)=(1-\varepsilon)P_0(W_k)$, $M_k=(1-\varepsilon)\E_0[xx^{\top}\one_{W_k}]$
and $v_k=(1-\varepsilon)\E_0[xy\one_{W_k}]$; the factors $1-\varepsilon$ cancel in
$\theta^{\infty}_k$, which proves the second claim of (iii).

\emph{Step 2: scales.} By the strong law, $n_k/n\to P(A_k)>0$, so $n_k\ge10$ for all large
$n$ and the pooled fallback is inactive. For $t\ge0$ put
$\hat F_k(t)=n_k^{-1}\#\{i:z_i=k,\,r_i\le t\}$ and $F_k(t)=P(r_\theta\le t\mid A_k)$; then
$\hat F_k(t)\to F_k(t)$ almost surely for each fixed $t$, as a ratio of two averages. Fix
$\delta>0$. By the uniqueness of the median, $F_k(m_k-\delta)<\tfrac12<F_k(m_k+\delta)$, so
for all large $n$, $\hat F_k(m_k-\delta)<\tfrac12<\hat F_k(m_k+\delta)$: fewer than half of
the residuals of group $k$ are at most $m_k-\delta$ and more than half are at most
$m_k+\delta$. Both middle order statistics, and hence $\hat m_k$, then lie in
$(m_k-\delta,m_k+\delta]$. Taking $\delta=1/j$, $j\ge1$, gives $\hat m_k\to m_k$ almost
surely, which is (i).

\emph{Step 3: flagged fraction.} Apply Lemma~\ref{lem:plugin} with $g\equiv1$, the
hypotheses holding by Step 1 and (i): $n^{-1}\#\{i:z_i=k,\,r_i\le c\hat s_k\}\to P(W_k)$ for
each $k$. Summing over $k$, $1-\hat\alpha_n(\theta)\to\sum_kP(W_k)=1-\alpha(\theta)$. By
Step 1, $\sum_kP(W_k)=(1-\varepsilon)\sum_kP_0(W_k)=(1-\varepsilon)(1-\kappa(\theta))$, which
rearranges to (ii).

\emph{Step 4: refit.} Apply Lemma~\ref{lem:plugin} with $g=xx^{\top}$ and with $g=xy$;
the integrability hypothesis holds by \eqref{eq:adp-cleanonly}, since
$\E_0\|x\|^{2}<\infty$ and $\E_0y^{2}<\infty$ under Assumption~\ref{ass:clean}. Thus
$n^{-1}\sum_ix_ix_i^{\top}\one\{z_i=k,\,i\notin\mathcal F_n\}\to M_k$ and
$n^{-1}\sum_ix_iy_i\one\{z_i=k,\,i\notin\mathcal F_n\}\to v_k$ almost surely. Since $M_k$ is
nonsingular, $P(W_k)>0$, so the number of unflagged units of group $k$ exceeds $d$ for all
large $n$, the fallback is inactive, and $\hat\theta_k$ is the image of the two averages
under the map $(M,v)\mapsto M^{-1}v$, which is continuous at $(M_k,v_k)$. This is (iii).
\end{proof}

Theorem~\ref{thm:stepfixed} identifies the limit of the step at any input satisfying the
assumptions. It does not say that the input of the step in Algorithm~\ref{alg:adaptive} is
close to $\theta^{0}$. The next theorem evaluates the limits at $\theta=\theta^{0}$, where they
can be made explicit, and quantifies the effect of overlap between the groups. Write
\[
\eta_k=P_0\bigl(h_{\theta^{0}}\ne k\bigm|z=k\bigr),\qquad
\eta=\sum_k\pi_k\eta_k=P_0\bigl(h_{\theta^{0}}\ne z\bigr),
\]
for the probability that a clean unit of group $k$, or a clean unit, is nearer to another
true surface than to its own. Let $p_k=P(A_k(\theta^{0}))$ and
\[
\delta_k=1-\frac{(1-\varepsilon)\pi_k(1-\eta_k)}{p_k},\qquad \nu_k=\delta_k+\tfrac12\eta_k,
\]
so that $\delta_k$ is the fraction of the units assigned to surface $k$ that are not clean
units of group $k$; it satisfies $\delta_kp_k\le(1-\varepsilon)\eta+\varepsilon$, since the
clean units of other groups assigned to $k$ have mass at most $(1-\varepsilon)\eta$. Finally
let
\[
q(u)=1.4826\,\Phi^{-1}\bigl(\tfrac{1+u}{2}\bigr),\qquad 0<u<1,
\]
an increasing function with $q(\tfrac12)=1.4826\,\Phi^{-1}(\tfrac34)=1.0000$ to four
decimals; $q(u)\sigma$ is the scale the step assigns to a $\mathcal N(0,\sigma^{2})$ residual
whose $u$-quantile in absolute value is used in place of the median.

\begin{theorem}[The step at the true parameter]
\label{thm:steptrue}
Let Assumption~\ref{ass:clean} hold, let $\theta=\theta^{0}$, and suppose that
$P_0(x^{\top}(\theta^{0}_k-\theta^{0}_j)=0)=0$ for $j\ne k$ (no two surfaces coincide on a
set of covariates of positive mass), that $\nu_k<\tfrac12$ for every $k$, and that
Assumption~\ref{ass:far} holds at $\theta^{0}$.
\begin{enumerate}[label=(\roman*),leftmargin=2.2em,itemsep=2pt]
\item For every $k$, $p_k>0$ and the median of $r_{\theta^{0}}$ on $A_k(\theta^{0})$ is
positive and unique, and
\[
q(\tfrac12-\nu_k)\;\le\;\frac{s_k(\theta^{0})}{\sigma_k}\;\le\;q(\tfrac12+\nu_k).
\]
\item The clean flagged fraction satisfies
\[
\sum_k\pi_k\,2\Phi\bigl(-c\,q(\tfrac12+\nu_k)\bigr)-\eta
\;\le\;\kappa(\theta^{0})\;\le\;
\sum_k\pi_k\,2\Phi\bigl(-c\,q(\tfrac12-\nu_k)\bigr)+\eta,
\]
and $\alpha(\theta^{0})=\varepsilon+(1-\varepsilon)\kappa(\theta^{0})$.
\item Suppose in addition that $\|x\|\le M_x$ $P_0$-almost surely and that
\[
\lambda_k:=\bigl\{2\Phi\bigl(c\,q(\tfrac12-\nu_k)\bigr)-1\bigr\}\,
\lambda_{\min}\bigl(\E_0[xx^{\top}\mid z=k]\bigr)-M_x^{2}\eta_k\;>\;0 .
\]
Then $M_k(\theta^{0})$ is nonsingular with $\lambda_{\min}(M_k(\theta^{0}))\ge(1-\varepsilon)\pi_k\lambda_k$,
Assumption~\ref{ass:reg} holds at $\theta^{0}$, Theorem~\ref{thm:stepfixed} applies, and
\[
\bigl\|\theta^{\infty}_k(\theta^{0})-\theta^{0}_k\bigr\|\;\le\;
\frac{2\,c\,s_k(\theta^{0})\,M_x}{\pi_k\lambda_k}\;\eta .
\]
\end{enumerate}
In particular, if $\varepsilon=0$, $\eta=0$ and every $\E_0[xx^{\top}\mid z=k]$ is nonsingular, as in (iii), then $s_k(\theta^{0})=q(\tfrac12)\sigma_k$,
$\kappa(\theta^{0})=2\Phi(-c\,q(\tfrac12))$ and $\theta^{\infty}(\theta^{0})=\theta^{0}$:
the true parameter is a fixed point of the population reweighting step, and the flagged
fraction is $2\Phi(-c)$ up to the rounding in $1.4826$.
\end{theorem}

\begin{proof}
Throughout, $h=h_{\theta^{0}}$, $r=r_{\theta^{0}}$, $A_k=A_k(\theta^{0})$,
$s_k=s_k(\theta^{0})$ and $G_k(t)=\Pr(\sigma_k|e|\le t)=2\Phi(t/\sigma_k)-1$.

\emph{Geometry of the assignment.} For $j\ne k$ put $\Delta_j(x)=x^{\top}(\theta^{0}_k-\theta^{0}_j)$,
which is nonzero $P_0$-almost surely. On $\{z=k\}$ we have $y-x^{\top}\theta^{0}_k=\sigma_ke$
and $y-x^{\top}\theta^{0}_j=\sigma_ke+\Delta_j$, and $|a|<|a+\Delta|$ if and only if
$\Delta(2a+\Delta)>0$. Hence, given $x$, the event $\{h=k\}$ is $\{-L(x)<\sigma_ke<U(x)\}$
with $L(x)=\min\{\Delta_j/2:\Delta_j>0\}$ and $U(x)=\min\{-\Delta_j/2:\Delta_j<0\}$, the
minimum over an empty set being $+\infty$; both are positive. Ties have probability zero.
On $\{z=k,h=k\}$ the residual is $r=\sigma_k|e|$.

\emph{Composition of $A_k$.} The mass of $A_k$ splits into the clean units of group $k$
with $h=k$, of mass $(1-\varepsilon)\pi_k(1-\eta_k)$, and the rest, of mass $\delta_kp_k$; the
rest consists of clean units of other groups, of mass at most
$(1-\varepsilon)\sum_j\pi_jP_0(h\ne j\mid z=j)=(1-\varepsilon)\eta$, and of contaminated
units. Since $\delta_k\le\nu_k<\tfrac12$, $p_k>0$ and $\eta_k<1$. The conditional distribution
function of $r$ on $A_k$ is therefore
\[
F_k(t)=(1-\delta_k)F^{\mathrm{own}}_k(t)+\delta_kF^{\mathrm{oth}}_k(t),\qquad
F^{\mathrm{own}}_k(t)=\frac{P_0(\sigma_k|e|\le t,\ h=k\mid z=k)}{1-\eta_k},
\]
with $F^{\mathrm{oth}}_k$ some distribution function. Since $e$ is independent of $(x,z)$,
$P_0(\sigma_k|e|\le t\mid z=k)=G_k(t)$, and $P_0(\sigma_k|e|\le t,\,h=k\mid z=k)$ lies between
$G_k(t)-\eta_k$ and $G_k(t)$. Hence
\begin{equation}
(1-\delta_k)\frac{G_k(t)-\eta_k}{1-\eta_k}\;\le\;F_k(t)\;\le\;
(1-\delta_k)\frac{G_k(t)}{1-\eta_k}+\delta_k .
\label{eq:adp-sandwich}
\end{equation}

\emph{(i) Uniqueness.} By the geometry above, $F^{\mathrm{own}}_k$ has density
\[
f^{\mathrm{own}}_k(t)=\frac{\varphi(t/\sigma_k)}{\sigma_k(1-\eta_k)}
\,\E_0\bigl[\one\{t<U(x)\}+\one\{t<L(x)\}\bigm|z=k\bigr],
\]
which is positive exactly when $P_0(\max\{L(x),U(x)\}>t\mid z=k)>0$; this property is
monotone in $t$, so $f^{\mathrm{own}}_k>0$ on $[0,T_k)$ and $f^{\mathrm{own}}_k=0$ on
$[T_k,\infty)$ for some $T_k\in(0,\infty]$. If the median of $F_k$ were not unique, $F_k$
would equal $\tfrac12$ on an interval of positive length. On $[0,T_k)$ the function $F_k$
is strictly increasing, so the interval lies in $[T_k,\infty)$, where
$F^{\mathrm{own}}_k=1$ and $F_k\ge1-\delta_k>\tfrac12$, a contradiction. The median is
positive because $F^{\mathrm{own}}_k(0)=0$ gives $F_k(0)\le\delta_k<\tfrac12$.

\emph{(i) Bounds.} Let $m_k$ be the median. If $F_k(t)\ge\tfrac12$ then, by the upper
bound in \eqref{eq:adp-sandwich}, $G_k(t)\ge u^{-}_k:=(\tfrac12-\delta_k)(1-\eta_k)/(1-\delta_k)$,
so $m_k\ge G_k^{-1}(u^{-}_k)$. If $G_k(t)\ge u^{+}_k:=\eta_k+(1-\eta_k)/(2(1-\delta_k))$ then,
by the lower bound, $F_k(t)\ge\tfrac12$, so $m_k\le G_k^{-1}(u^{+}_k)$; here $0<u^{-}_k$ and
$u^{+}_k<1$ because $\delta_k<\tfrac12$ and $\eta_k<1$. Now
$(\tfrac12-\delta_k)/(1-\delta_k)\ge\tfrac12-\delta_k$ and
$1/(2(1-\delta_k))\le\tfrac12+\delta_k$ for $\delta_k\le\tfrac12$, whence
$u^{-}_k\ge(\tfrac12-\delta_k)(1-\eta_k)\ge\tfrac12-\nu_k>0$ and
$u^{+}_k\le\eta_k+(1-\eta_k)(\tfrac12+\delta_k)\le\tfrac12+\nu_k<1$. Since
$G_k^{-1}(u)=\sigma_k\Phi^{-1}((1+u)/2)$ is increasing and $s_k=1.4826\,m_k$, the bounds
follow.

\emph{(ii).} Split $\kappa(\theta^{0})=P_0(r>cs_h)$ according to whether $h=z$. On
$\{h=z=k\}$, $r=\sigma_k|e|$ and $s_h=s_k$, so
$P_0(h=z,\ r>cs_h)=\sum_k\pi_kP_0(\sigma_k|e|>cs_k,\ h=k\mid z=k)$, which lies between
$\sum_k\pi_k\{G'_k-\eta_k\}$ and $\sum_k\pi_kG'_k$ with
$G'_k=\Pr(\sigma_k|e|>cs_k)=2\Phi(-cs_k/\sigma_k)$. The event $\{h\ne z\}$ contributes
between $0$ and $\eta$. Inserting the bounds of (i) on $s_k/\sigma_k$ into $G'_k$ gives
(ii); the expression for $\alpha(\theta^{0})$ is Theorem~\ref{thm:stepfixed}(ii), whose
Assumption~\ref{ass:reg} is verified in (iii) below, or directly from
\eqref{eq:adp-cleanonly}.

\emph{(iii) Nonsingularity.} By \eqref{eq:adp-cleanonly},
$M_k=(1-\varepsilon)\E_0[xx^{\top}\one_{W_k}]$, and $W_k\supseteq\{z=k,\,h=k,\,\sigma_k|e|\le cs_k\}$.
In the positive semidefinite order,
\begin{multline*}
\E_0\bigl[xx^{\top}\one\{z=k,\,h=k,\,\sigma_k|e|\le cs_k\}\bigr]\\
=\E_0\bigl[xx^{\top}\one\{z=k,\,\sigma_k|e|\le cs_k\}\bigr]
-\E_0\bigl[xx^{\top}\one\{z=k,\,h\ne k,\,\sigma_k|e|\le cs_k\}\bigr],
\end{multline*}
where the first term equals $\pi_kG_k(cs_k)\E_0[xx^{\top}\mid z=k]$ by the independence of
$e$ and $(x,z)$, and the second is at most $M_x^{2}\pi_k\eta_kI$ because
$xx^{\top}\preceq\|x\|^{2}I$. With $G_k(cs_k)=2\Phi(cs_k/\sigma_k)-1\ge2\Phi(cq(\tfrac12-\nu_k))-1$
by (i), $\lambda_{\min}(M_k)\ge(1-\varepsilon)\pi_k\lambda_k>0$. The other two parts of
Assumption~\ref{ass:reg} were shown in (i).

\emph{(iii) Bias.} $\theta^{\infty}_k-\theta^{0}_k=M_k^{-1}\E[x(y-x^{\top}\theta^{0}_k)\one_{W_k}]$,
and by \eqref{eq:adp-cleanonly} the expectation is $(1-\varepsilon)$ times its $P_0$-version.
Split the latter by $z$. On $\{z=k\}$, $y-x^{\top}\theta^{0}_k=\sigma_ke$ and
$W_k\cap\{z=k\}=\{z=k,\,h=k,\,|e|\le cs_k/\sigma_k\}$, so
\begin{multline*}
\E_0\bigl[xe\,\one\{z=k,\,h=k,\,|e|\le cs_k/\sigma_k\}\bigr]\\
=\E_0\bigl[x\one\{z=k\}\bigr]\,\E\bigl[e\one\{|e|\le cs_k/\sigma_k\}\bigr]
-\E_0\bigl[xe\,\one\{z=k,\,h\ne k,\,|e|\le cs_k/\sigma_k\}\bigr].
\end{multline*}
The first term vanishes by the symmetry of $e$; the second has norm at most
$(cs_k/\sigma_k)M_x\pi_k\eta_k$, so the contribution of $\{z=k\}$ is at most
$cs_kM_x\pi_k\eta_k$ in norm. On $\{z=j\}$, $j\ne k$, the factor $|y-x^{\top}\theta^{0}_k|=r\le cs_k$
on $W_k$, so the contribution is at most $cs_kM_xP_0(z=j,\,h=k)$ in norm, and summing over
$j\ne k$ gives at most $cs_kM_x\eta$. Altogether
$\|\E[x(y-x^{\top}\theta^{0}_k)\one_{W_k}]\|\le(1-\varepsilon)cs_kM_x(\pi_k\eta_k+\eta)\le2(1-\varepsilon)cs_kM_x\eta$,
and dividing by $\lambda_{\min}(M_k)\ge(1-\varepsilon)\pi_k\lambda_k$ gives the bound.

\emph{The special case.} If $\varepsilon=0$ and $\eta=0$ then $\delta_k=\nu_k=0$, so
$s_k=q(\tfrac12)\sigma_k$ by (i), $\kappa=2\Phi(-cq(\tfrac12))$ by (ii), and the bias bound
is zero.
\end{proof}

Three comments. First, the constant in the limit of the flagged fraction is
$2\Phi(-c\,s_k/\sigma_k)$ and not $2\Phi(-c)$. On the units assigned to surface $k$ the
minimum-over-surfaces residual is the residual from surface $k$ itself, so the per-group
median is a median of own-surface residuals; it is pulled down because the clean units of
group $k$ with the largest errors are nearer to another surface and are assigned there, and
pulled up by the units of other groups and by the contamination assigned to surface $k$.
Theorem~\ref{thm:steptrue}(i) bounds the net effect by $\nu_k=\delta_k+\eta_k/2$, which is
of the order of $\eta+\varepsilon$. The first effect is intrinsic to any scale computed
within an assigned group and vanishes only when the groups do not overlap. Second, the refit
is exactly unbiased at $\theta^{0}$ when the groups do not overlap and moves by $O(\eta)$
when they do; the source is the asymmetric truncation of the error on $\{h=k\}$, not the
contamination, which the step removes entirely under Assumption~\ref{ass:far}. Third,
Algorithm~\ref{alg:adaptive} reports the flagged fraction recomputed at the refitted
parameter, that is $\hat\alpha_n(\hat\theta)$ rather than $\hat\alpha_n(\theta)$;
Theorem~\ref{thm:stepfixed} concerns the flag set the refit uses. The limit of the reported
fraction is $\alpha(\theta^{\infty}(\theta))$ whenever $\alpha$ is continuous at
$\theta^{\infty}(\theta)$ and the convergence of $\hat\alpha_n$ is uniform near it, which we
do not prove.


\emph{Scope of Proposition~\ref{thm:step} of the paper.} The proposition is a law of large numbers for one application of the map at a fixed input. It
says what the flagged fraction estimates under Gaussian errors, $\varepsilon$ plus the fraction
$\kappa$ of clean units flagged, which is close to $2\Phi(-c)\approx0.012$ when the groups
overlap little, and that the contamination is removed entirely once it lies beyond the
cut-off, leaving a bias of order $\eta$ from the overlap of the groups. Two restrictions of
Theorem~\ref{thm:steptrue} should be stated. First, the constant is
$\nu_k=\delta_k+\eta_k/2$, where $\delta_k$ is the fraction of the units assigned to surface
$k$ that are not clean units of group $k$ and $\eta_k$ the overlap of group $k$; it satisfies
$\nu_k\le\{(1-\varepsilon)\eta+\varepsilon\}/p_k+\eta_k/2$ with $p_k$ the mass assigned to
surface $k$, so it is of order $\eta+\varepsilon$ divided by $p_k$ and is not small for a
small group. Second, the theorem requires $\nu_k<\tfrac12$ for every group, a breakdown
condition per group, and this fails for a small group with moderate contamination: with
$\pi_k=0.15$, $\varepsilon=0.20$ and the contamination split evenly between the two assignment
regions of two groups, $\delta_k$ is about $0.45$ without overlap and $\nu_k$ exceeds
$\tfrac12$ once about a tenth of the group's clean units lie nearer to the other line. The
theorem does not cover that regime. Three further things it does not cover should be stated. It does not show that ESF delivers
an input close to $\theta^{0}$. The code applies the map more than once: once at the end of
ESF, twice after TCLUST-REG in the reweighted rows of Section~4 of the paper, and three
times to each candidate of the recovery step. And the reported flagged fraction $\hat\alpha$
is evaluated at the refitted parameter $\hat\theta$, whereas the proposition concerns the
flags computed at the input $\theta$; the two limits agree under a continuity condition that
we do not prove (the third comment above).

\subsection*{S2.3 The group-recovery step}
\label{app:recovery}

Proposition~\ref{prop:recover-value} of the paper is for fixed parameters; the step
reads them off the data, and the reweighting refit of the candidate is not covered. The local
peak test guards against a slice of a band of gross outliers, which the likelihood would value
like a group; for a residual density fixed in advance it is a concave maximisation whose
solution is zero for a band that fills the window and the Gaussian share for a Gaussian group
over a uniform background, as follows. Let $W=3cs$, $u=1/(2W)$, $g_W$ the $\mathcal N(0,s^{2})$ density
truncated to $[-W,W]$, and for a residual density $p$ on the window let $w^\ast(p)$ maximise
$L_p(w)=\int p\log\{wg_W+(1-w)u\}$ over $[0,1]$.

\begin{proposition}
\label{prop:recover-peak}
\begin{enumerate}[label=(\roman*),leftmargin=2.2em,itemsep=2pt]
\item $L_p$ is strictly concave, so $w^\ast(p)$ is unique, and the EM iteration in $w$
started at $\tfrac12$ converges monotonically to it; the same holds for the sample version,
whose maximiser converges to $w^\ast(p)$ almost surely for independent residuals of density $p$.
\item If $|r|$ is stochastically smaller under $\tilde p$ than under $p$, then
$w^\ast(\tilde p)\ge w^\ast(p)$.
\item If $p=u$, then $w^\ast(p)=0$; if $p$ is uniform on $[-V,V]$, $V\le W$, then $w^\ast$
is nonincreasing in $V$.
\item If $p=w_0g_W+(1-w_0)u$, then $w^\ast(p)=w_0$.
\end{enumerate}
\end{proposition}

The test thus passes a group whose units are at least half of those in the window and fails
a slice of a band that fills it. Between the two, solving $L_p'(w)=0$ for $c=2.5$, a uniform
band of half-width $V$ has $w^\ast\ge\tfrac12$ exactly when $V\le3.44\,s$: a band narrower
than about seven noise scales passes the test and is left to the scale restriction and the
likelihood. The implementation uses the untruncated normal density, which differs from
$g_W$ by the mass $2\Phi(-3c)\approx6\times10^{-14}$ outside the window; (ii) and (iii) hold
for it as stated. The proposition concerns a given residual density: the effect of the RANSAC
selection on $p$ is not analysed, and the fifty EM iterations of the implementation
approximate the limit in (i).


\emph{Proposition~\ref{prop:recover-value} of the paper (precise statement).} Let the units have law $P$.
\begin{enumerate}[label=(\roman*),leftmargin=2.2em,itemsep=2pt]
\item Let $G$ be an event with $P(G)=\pi_\ast$, on which $y=x^{\top}\theta_\ast+\sigma_\ast e$
with $e\sim\mathcal N(0,1)$ independent of $x$ given $G$. Let $F$ have noise weight
$\pi_0>\pi_\ast$, and let $F'$ be $F$ with the line $(\theta_\ast,\sigma_\ast,\pi_\ast)$
added and the noise weight lowered to $\pi_0-\pi_\ast$. Then \eqref{eq:recover-gain} of the paper holds
with $\xi=\E[\log\{(R/\pi_0)f_F(u)\}\mid G]\ge0$, the extent to which $F$ already explains
$G$, and $\nu=\E[(\pi_0/R)f_F(u)^{-1}\one_{G^{c}}]\le1-\pi_\ast$, the noise mass that $F$
attributes to the units outside $G$; a bound on $\xi$ in terms of the distance of $G$ from
the lines of $F$ is given below.
\item Let $S$ be an event with $P(S)=\pi_S$, on which $y=x^{\top}\theta_S+\sigma_Se$ as in
(i). Let $F$ contain the lines $\theta_S\pm\tfrac{\delta}{2}e_1$, $e_1$ the intercept
coordinate, each with scale $\sigma_S$ and proportion $\pi_S/2$, and let $F''$ replace them
by the line $(\theta_S,\sigma_S,\pi_S)$, all else equal. With $a=\delta/(2\sigma_S)$,
$-\varphi(1)a^{2}\pi_S\le L(F'')-L(F)\le(\varphi(1)a^{2}+\tfrac14a^{4})\pi_S$. If the density
$r$ of the other components on $S$ is replaced by $tr$ and $t\downarrow0$, the difference
tends to $\pi_S$ times the Kullback--Leibler divergence from $\mathcal N(0,1)$ to
$\tfrac12\mathcal N(-a,1)+\tfrac12\mathcal N(a,1)$, which lies in $[0,a^{4}/4]$.
\end{enumerate}
Hence, if a candidate $F_1$ arises from the original $F_0$ by the two
operations and the right side of \eqref{eq:recover-gain} exceeds $\varphi(1)a^{2}\pi_S$, then
almost surely $\ell_n(F_1)-\ell_n(F_0)>\tfrac{d+1}{2}\log n$ for all large $n$, and the
candidate is accepted if its scales are admissible.

\emph{Proposition~\ref{prop:recover-value}(i).} Write $f=f_F$ and $f'=f_{F'}$, so that
$f'=f-\pi_\ast/R+(\pi_\ast/\sigma_\ast)\varphi((y-x^{\top}\theta_\ast)/\sigma_\ast)$, and
split $L(F')-L(F)=\E[(\log f'-\log f)\one_G]+\E[(\log f'-\log f)\one_{G^{c}}]$.

On $G$, $y-x^{\top}\theta_\ast=\sigma_\ast e$, and dropping the lines of $F$ from $f'$,
\[
f'\;\ge\;\frac{\pi_\ast}{\sigma_\ast}\varphi(e)+\frac{\pi_0-\pi_\ast}{R}
=\frac{\pi_\ast}{\sigma_\ast}\varphi(e)\bigl(1+\tau e^{e^{2}/2}\bigr),
\qquad \tau=\frac{(\pi_0-\pi_\ast)\sigma_\ast\sqrt{2\pi}}{\pi_\ast R}.
\]
Since $\E\log\varphi(e)=-\tfrac12\log(2\pi)-\tfrac12\E e^{2}=-\tfrac12\log(2\pi e)$,
\[
\E[\log f'\mid G]\;\ge\;\log\frac{\pi_\ast}{\sigma_\ast}-\tfrac12\log(2\pi e)+\rho,
\qquad \rho=\E\log\bigl(1+\tau e^{e^{2}/2}\bigr)\ge0 .
\]
For $v\ge0$, $\log(1+v)\le2\sqrt v$ (the difference vanishes at $0$ and has derivative
$v^{-1/2}-(1+v)^{-1}>0$), and $\E e^{e^{2}/4}=(1-\tfrac12)^{-1/2}=\sqrt2$; hence
$\rho\le2\sqrt\tau\,\E e^{e^{2}/4}=2\sqrt{2\tau}$. Next, $f=(\pi_0/R)\{1+(R/\pi_0)\sum_k
(\pi_k/s_k)\varphi((y-x^{\top}\theta_k)/s_k)\}$, so $\E[\log f\mid G]=\log(\pi_0/R)+\xi$. By
Jensen's inequality, $\xi\le\log\{1+(R/\pi_0)\sum_k\pi_k\E[s_k^{-1}\varphi((y-x^{\top}\theta_k)/s_k)\mid G]\}$,
and on $G$, $y-x^{\top}\theta_k=x^{\top}(\theta_\ast-\theta_k)+\sigma_\ast e$ with $e$
independent of $x$; integrating over $e$ first, the convolution of two centred normal
densities gives $\E[s_k^{-1}\varphi((y-x^{\top}\theta_k)/s_k)\mid x,G]=\varphi_k(x^{\top}(\theta_\ast-\theta_k))$,
where $\varphi_k$ is the density of $N(0,s_k^{2}+\sigma_\ast^{2})$. Hence
\[
\xi\le\log\Bigl\{1+\frac{R}{\pi_0}\sum_{k}\pi_k\,\E\bigl[\varphi_k\bigl(x^{\top}(\theta_\ast-\theta_k)\bigr)\bigm| G\bigr]\Bigr\},
\]
the bound on $\xi$ referred to in Proposition~3(i) of the paper: it is small when the lines of
$F$ are many scales $(s_k^2+\sigma_\ast^2)^{1/2}$ away from $G$. Multiplying by $P(G)=\pi_\ast$, the first expectation is at
least $\pi_\ast$ times the brace in \eqref{eq:recover-gain} plus $\pi_\ast\rho$; the
statement drops $\rho\ge0$, and the bound $\rho\le2\sqrt{2\tau}$ shows that the leading
term is exact up to $\xi$ and the two dropped nonnegative terms, $\rho$ and the density
of the lines of $F$ on $G$ in $f'$.

On $G^{c}$, $f'\ge f-\pi_\ast/R>0$ because $f\ge\pi_0/R>\pi_\ast/R$, so
$\log f'-\log f\ge-\log\{f/(f-\pi_\ast/R)\}$ and the second expectation is at least
$-\Lambda$, $\Lambda=\E[\log\{f/(f-\pi_\ast/R)\}\one_{G^{c}}]$. With $t=(\pi_0/R)/f\in(0,1]$
and $b=\pi_\ast/\pi_0<1$, $\log\{f/(f-\pi_\ast/R)\}=-\log(1-bt)$; the function
$t\mapsto-\log(1-bt)$ is convex on $[0,1]$ and vanishes at $0$, so it is at most $t$ times
its value at $1$, namely $t\log\{\pi_0/(\pi_0-\pi_\ast)\}$. Taking expectations over $G^{c}$
gives $\Lambda\le\nu\log\{\pi_0/(\pi_0-\pi_\ast)\}$, and $\nu\le P(G^{c})$ since $t\le1$. \qed

\emph{Proposition~\ref{prop:recover-value}(ii).} On $S$, the residuals from the two lines are
$\sigma_Se\mp\delta/2$, so their joint density is
\[
A=\frac{\pi_S}{2\sigma_S}\bigl\{\varphi(e-a)+\varphi(e+a)\bigr\}
=\frac{\pi_S}{\sigma_S}\varphi(e)\,e^{-a^{2}/2}\cosh(ae),
\qquad\text{against}\qquad B=\frac{\pi_S}{\sigma_S}\varphi(e)
\]
for the single line, and $\log(B/A)=a^{2}/2-\log\cosh(ae)$. Let $r\ge0$ be the density of
the other components, the same in both fits. Off $S$ the two densities coincide, and on $S$,
$\log f_{F''}-\log f_F=\psi$ with $\psi=\log\{(B+r)/(A+r)\}$. Since
$t\mapsto\log\{(B+t)/(A+t)\}$ is monotone on $[0,\infty)$ with limit $0$, $\psi$ lies between
$\min\{\log(B/A),0\}$ and $\max\{\log(B/A),0\}$. Hence
\[
-\pi_SD_-(a)\;\le\;L(F'')-L(F)=\pi_S\E[\psi\mid S]\;\le\;\pi_SD_+(a),
\]
where $D_\pm(a)=\E[(\pm\{\tfrac12a^{2}-\log\cosh(ae)\})^{+}]$.
Two inequalities for $x\in\R$: $\log\cosh x\le x^{2}/2$, because both sides vanish at $0$ and
$\tanh x\le x$ for $x\ge0$; and $\log\cosh x\ge x^{2}/2-x^{4}/12$, because the difference
vanishes at $0$ and its derivative $\tanh x-x+x^{3}/3$ is nondecreasing on $[0,\infty)$, its
own derivative being $x^{2}-\tanh^{2}x\ge0$. By the first, $D_-(a)\le\E[(a^{2}e^{2}/2-a^{2}/2)^{+}]
=\tfrac12a^{2}\E(e^{2}-1)^{+}=\varphi(1)a^{2}$, since
$\E(e^{2}-1)^{+}=2\int_1^\infty(z^{2}-1)\varphi(z)\,dz=2\varphi(1)$. Further
$D_+(a)-D_-(a)=\E\log(B/A)=\E\log\{\varphi(e)/\tfrac12(\varphi(e-a)+\varphi(e+a))\}$ is the
Kullback--Leibler divergence stated, hence nonnegative, and by the second inequality it is at
most $\E[a^{4}e^{4}]/12=a^{4}/4$. So $D_+(a)\le\varphi(1)a^{2}+a^{4}/4$, which gives the two
bounds. Finally, if $r$ is replaced by $tr$, then $\psi\to\log(B/A)$ pointwise as
$t\downarrow0$, dominated by $|\log(B/A)|\le a^{2}e^{2}/2+a^{2}/2$, which is integrable; by
dominated convergence the difference tends to $\pi_S\E\log(B/A)$. \qed

\emph{Acceptance.} If $F_1$ arises from $F_0$ by (i) and then (ii), then
$L(F_1)-L(F_0)$ is at least the right side of \eqref{eq:recover-gain} minus
$\varphi(1)a^{2}\pi_S$, which is positive by assumption. By the strong law, applied to the
bounded summands $\log f_{F_1}(u_i)$ and $\log f_{F_0}(u_i)$,
$\{\ell_n(F_1)-\ell_n(F_0)\}/n\to L(F_1)-L(F_0)>0$ almost surely, while
$\tfrac{d+1}{2}\log n=o(n)$. \qed

\emph{Proposition~\ref{prop:recover-peak}.} On the window, $g_W$ lies between
$\varphi(3c)/(s\kappa)$ and $\varphi(0)/(s\kappa)$ with $\kappa=1-2\Phi(-3c)$, and $u>0$, so
$\log\{wg_W+(1-w)u\}$ is bounded on $[-W,W]\times[0,1]$ and $L_p$ is finite. Write
$q_w=wg_W+(1-w)u$.

\emph{(i).} For fixed $r$, $w\mapsto\log q_w(r)$ is concave, and strictly so unless
$g_W(r)=u$, which happens for at most two values of $r$, a $p$-null set. Hence $L_p$ is
strictly concave and has a unique maximiser. Its derivative
$L_p'(w)=\int p\,(g_W-u)/q_w$ is continuous and strictly decreasing, so
$w^\ast=\sup\{w\in[0,1]:L_p'(w)\ge0\}$, the supremum of the empty set being $0$. The EM map
is $M(w)=\int p\,wg_W/q_w$, and
\[
M(w)-w=\int p\,\frac{wg_W-wq_w}{q_w}=w(1-w)L_p'(w),
\qquad
M'(w)=\int p\,\frac{g_Wu}{q_w^{2}}>0 .
\]
Thus $M$ is increasing, its fixed points are $0$, $1$ and the zero of $L_p'$ in $(0,1)$
when there is one, and $M(w)>w$ exactly when $L_p'(w)>0$, that is when $w<w^\ast$. Start at
$w_0=\tfrac12$. If $w_0<w^\ast$, then $w_0<M(w_0)\le M(w^\ast)=w^\ast$ when $w^\ast<1$,
and $M(w_0)\le1$ otherwise; by induction the iterates increase and stay at most $w^\ast$, so
they converge to a fixed point in $(\tfrac12,w^\ast]$, which is $w^\ast$. If $w_0>w^\ast$,
they decrease and stay at least $w^\ast$, and converge to a fixed point in $[w^\ast,\tfrac12)$,
which is $w^\ast$ when $w^\ast>0$ and $0$ when $w^\ast=0$. The sample criterion
$L_N(w)=N^{-1}\sum_{i\le N}\log q_w(r_i)$ has the same properties, with $p$ replaced by the
empirical law, as soon as one $r_i$ has $g_W(r_i)\ne u$, which is almost sure when the
$r_i$ have a density; the implemented iteration is its EM map, so the same argument applies
to its maximiser $\hat w_N$. If the $r_i$ are independent with density $p$, then
$L_N(w)\to L_p(w)$ almost surely for each $w$ in a countable dense subset of $[0,1]$, by the
strong law for bounded summands; the functions $w\mapsto\log q_w(r)$ have derivative bounded
by $\max(g_W,u)/\min(g_W,u)$ uniformly in $r$, so the $L_N$ are equi-Lipschitz and the
convergence is uniform on $[0,1]$; uniform convergence to a continuous function with a unique
maximiser on a compact set gives $\hat w_N\to w^\ast(p)$.

\emph{(ii).} For $t\ge0$ put $h_w(t)=(g_W(t)-u)/(wg_W(t)+(1-w)u)$. As a function of
$g\in(0,\infty)$, $(g-u)/(wg+(1-w)u)$ has derivative $u/(wg+(1-w)u)^{2}>0$, and $g_W(t)$ is
decreasing in $t$, so $h_w$ is decreasing and bounded. Since
$L_p'(w)=\E_p[h_w(|r|)]$, the stochastic order gives $L_{\tilde p}'(w)\ge L_p'(w)$ for every
$w$, hence $\{w:L_{\tilde p}'(w)\ge0\}\supseteq\{w:L_p'(w)\ge0\}$ and
$w^\ast(\tilde p)\ge w^\ast(p)$ by the characterisation in (i).

\emph{(iii).} If $p=u$, then $L_p'(0)=\int u\,(g_W-u)/u=\int_{-W}^{W}g_W-1=0$, so
$L_p'(w)<0$ for $w>0$ and $w^\ast=0$. With the untruncated normal density $g$ in place of
$g_W$ the same computation gives $L_p'(0)=\kappa-1<0$, and again $w^\ast=0$. If $p$ is uniform
on $[-V,V]$, then $|r|$ is stochastically smaller for smaller $V$, and (ii) applies; the
proof of (ii) uses only that the peak density is decreasing in $|r|$, so it holds with $g$
in place of $g_W$ as well.

\emph{(iv).} If $p=w_0g_W+(1-w_0)u=q_{w_0}$, then $L_p'(w_0)=\int(g_W-u)=0$, so $w_0$ is the
zero of $L_p'$ and $w^\ast=w_0$. \qed

\emph{Numerical values.} The constants quoted in Section~\ref{sec:recovery-theory} were
obtained by solving $L_p'(w)=0$ by quadrature and bisection with $s=1$ and $c=2.5$. For the
uniform band of half-width $V$, $w^\ast$ is $0$, $0.08$, $0.18$, $0.34$, $0.71$ and $1$ at
$V=7.5$, $6$, $5$, $4$, $3$ and $2.5$, and equals $\tfrac12$ at $V=3.44$. For a slice whose
density on $|r|\le cs$ exceeds that on the rest of the window by a factor $1+\beta$, as the
RANSAC selection tends to produce, $w^\ast$ reaches $\tfrac12$ at $\beta=5.2$, that is when
the central third of the window holds $76\%$ of its units against $33\%$ for a uniform slice.
For a mixture of a Gaussian share $w_0\in\{0.3,0.5,0.7\}$ and a uniform background the
computed maximiser equals $w_0$ to six decimals, with $g_W$ and with the untruncated density
alike. The bounds of Proposition~\ref{prop:recover-value} were checked by quadrature and by
Monte Carlo with $2\times10^{6}$ units: with $R/\sigma_\ast=30$, $\pi_\ast=0.1$,
$\pi_0=0.15$ and two lines seven scales away, the mean gain per unit of $G$ was $1.6695$
against a leading term of $1.5768$, $\rho=0.0928$ and $\xi\le1.2\times10^{-4}$; with the
lines three scales away the gain fell to $0.90$ and the Jensen bound on $\xi$ was $1.85$
against a true value of $1.47$. For the un-splitting, $D_+-D_-$ and $D_-$ were within their
bounds at $a\in\{0.1,0.25,0.5,1\}$, and with a noise floor present the mean change per unit
of $S$ at $a=0.5$ was $0.016$, inside $[-0.049,0.061]$, against $0.012$ for the divergence.


